% ===================================================================
%  तक्ता क्रमांक २ — मानवी शरीर. — मधली सुट्टी. खेळ.
%
%  Traduction, faite en 2026, en marathi d'aujourd'hui, sur l'ido de
%  texto/io. Regles de la colonne : voir 10-takta-01.tex.
%
%  LE CORPS HUMAIN EST LE TABLEAU OU LE MARATHI SE SEPARE LE PLUS DU
%  HINDI, ET AUCUN DE CES MOTS N'EST UN EMPRUNT : ce sont ceux que
%  tout le monde emploie, et le hindi en emploie d'autres.
%
%    डोके        la tete      हिन्दी : सिर
%    केस         les cheveux  हिन्दी : बाल
%    कपाळ        le front     हिन्दी : माथा
%    डोळा        l'œil        हिन्दी : आँख
%    हनुवटी      le menton    हिन्दी : ठोड़ी
%    मिशी        la moustache हिन्दी : मूँछ
%    कल्ले       les favoris
%    ओठ          les levres   हिन्दी : होंठ
%    मान         le cou       हिन्दी : गर्दन
%    घसा         la gorge     हिन्दी : गला
%    मानगूट      la nuque
%    फासळी       la cote      हिन्दी : पसली
%    पाठ         le dos       हिन्दी : पीठ
%    कोपर        le coude     हिन्दी : कोहनी
%    मनगट        le poignet   हिन्दी : कलाई
%    गुडघा       le genou     हिन्दी : घुटना
%    पोटरी       le mollet    हिन्दी : पिंडली
%    कातडी       la peau      हिन्दी : चमड़ा
%    पाऊल        le pied      हिन्दी : पैर
%    तळपाय       la plante    हिन्दी : तलवा
%    टाच         le talon     हिन्दी : एड़ी
%    मेंदू        le cerveau   हिन्दी : दिमाग
%
%  C'EST LE MEME RESULTAT QU'AU TABLEAU 2 EGYPTIEN, ou vingt-deux mots
%  du corps separaient l'egyptien de l'arabe standard. La planche du
%  corps est l'endroit ou une langue se prouve : on n'y choisit pas
%  son registre, on y dit ce qu'on dit.
%
%  ET LES JEUX PORTENT LEURS NOMS MARATHIS :
%
%    चेंडू (17, 22)  la balle    हिन्दी : गेंद
%    भोवरा (18)     la toupie
%    गोट्या (21)    les billes
%    झोपाळा (35)    la balancoire
%    पतंग (33)      le cerf-volant
%    फुगा (41)      le ballon
%    पोहणे (40)     la nage     हिन्दी : तैरना
%    शिडी (2)       l'echelle   हिन्दी : सीढ़ी
%
%  ET DEUX JEUX QUE LE MARATHI DISTINGUE EXACTEMENT COMME L'IDO :
%  « celo-trovo » (25) est लपाछपी, se cacher et se chercher ;
%  « blindo-ludo » (31) est आंधळी कोशिंबीर — « la salade aveugle » —
%  les yeux bandes. Le francais de 1926 dit cache-cache et
%  colin-maillard, et le marathi a ses deux noms au meme endroit,
%  dont l'un est une image que personne ne remarque plus. C'est la
%  troisieme colonne de suite ou la paire tombe juste : le cantonais
%  et l'egyptien l'avaient donnee aussi.
%
%  ET UN TROU, COMME LE FLEAU DU TABLEAU 8 EGYPTIEN. Le « baro-ludo »
%  (38) — le jeu de barres, deux camps qui se poursuivent — n'a pas de
%  nom marathi, parce que le Maharashtra joue a खो-खो et a
%  आट्यापाट्या, qui n'en sont pas. Les nommer ainsi aurait change la
%  CHOSE ; on ecrit donc « बार-खेळ », qui decrit au lieu de nommer.
% ===================================================================

%%K t02-apar-1 apar
\VUsaut{4.0mm}
\VUtitre{15.0pt}{60}{तक्ता क्रमांक २}
\VUsaut{2.0mm}
\VUtitre{11.4pt}{0}{मानवी शरीर. --- मधली सुट्टी.}
\VUtitre{11.4pt}{0}{खेळ.}

%%K t02-tit-0 sub
\VUtitre{13.2pt}{0}{मानवी शरीर.}
\VUsaut{2.4mm}
\VUcentre{10.2pt}{0}{\textit{(निसर्गविज्ञानाचा साधा धडा.)}}
\VUsaut{3.0mm}

%%K t02-tit-1 sub pos=t02-01-1
\VUpk{10.2pt}{40}{डोके.}
\VUsaut{3.0mm}

%%K t02-01-1 p
१. --- मानवी शरीराचे मुख्य भाग म्हणजे: \VUgras{डोके}, \VUgras{धड} आणि
\VUgras{अवयव}.

%%K t02-02-1 p
२. --- डोक्याचे दोन भाग असतात: \VUgras{कवटी}, जिच्यात \VUgras{मेंदू} असतो आणि
जिच्यावर \VUgras{केस} असतात, आणि \VUgras{चेहरा}.

%%K t02-03-1 p
३. --- आपल्या चित्रात कवटीही दिसत नाही आणि मेंदूही दिसत नाही; पण चेहऱ्याचे
महत्त्वाचे भाग अगदी स्पष्ट दिसतात.

%%K t02-04-1 p
४. --- सर्वांत आधी \VUgras{कपाळ}, जे प्रबळ बुद्धीचे आणि सतत चालणाऱ्या विचाराचे
चिन्ह आहे. \VUgras{कानशिले} अगदी मोकळी आहेत. \VUgras{डोळा} तेजस्वी आणि पुरता
उघडलेला आहे. दाट \VUgras{भुवई} आणि लांब \VUgras{पापण्या} त्याचे रक्षण करतात.

%%K t02-05-1 p
५. --- मग \VUgras{नाक} आणि \VUgras{नाकपुड्या}. नाकाचा उपयोग वास घेण्यासाठी आणि
श्वास घेण्यासाठी होतो. नाकाच्या दोन्ही बाजूंना \VUgras{गाल} आहेत, आणखी पलीकडे
\VUgras{कान}. चेहऱ्याच्या खालच्या भागात \VUgras{तोंड} आणि \VUgras{हनुवटी} आहेत.

%%K t02-06-1 p
६. --- ते आपल्याला फारसे नीट दिसत नाहीत, कारण \VUgras{मिशी} आणि \VUgras{कल्ले}
ते झाकून टाकतात. पण आपल्याला माहीत आहे की तोंडात दोन \VUgras{ओठ} असतात,
\VUgras{वरचा जबडा} आणि \VUgras{खालचा जबडा} त्यांच्या \VUgras{दातांसह}, आणि
\VUgras{जीभ}. तोंडाने आपण बोलतो आणि खातो. जिभेने आपण अन्नाची चव घेतो.

%%K t02-07-1 p
७. --- डोळा, कान, नाक आणि तोंड, बोटांप्रमाणेच, पाच ज्ञानेंद्रियांची इंद्रिये
आहेत: \VUgras{दृष्टी}, \VUgras{श्रवण}, \VUgras{घ्राण}, \VUgras{चव},
\VUgras{स्पर्श}.

%%K t02-08-1 p
८. --- म्हणून डोके हा मानवी शरीराच्या सर्वांत रोचक भागांपैकी एक आहे.

%%K t02-tit-2 sub
\VUsaut{4.4mm}
\VUpk{10.2pt}{40}{धड.}
\VUsaut{3.0mm}

%%K t02-09-1 p
९. --- \VUgras{मान} डोक्याला धडाशी जोडते. तिच्यात \VUgras{घसा} आणि
\VUgras{मानगूट} येतात.

%%K t02-10-1 p
१०. --- धडातही पुष्कळ महत्त्वाचे अवयव असतात. \VUgras{छातीत} \VUgras{फुप्फुसे}
असतात, ज्यांनी आपण श्वास घेतो, आणि \VUgras{हृदय} असते, जे जीवनाचे केंद्र आहे.
हृदयाचे ठोके थांबले की सजीव मरतो.

%%K t02-11-1 p
११. --- \VUgras{फासळ्या} छातीला आधार देतात.

%%K t02-12-1 p
१२. --- धडाचे उरलेले भाग म्हणजे \VUgras{कूस}, \VUgras{पाठ}, \VUgras{खांदे} आणि
\VUgras{पोट}. झोपण्यासाठी आपण कुशीवर किंवा पाठीवर आडवे होतो, आणि ओझी आपण
खांद्यावर वाहतो.

%%K t02-13-1 p
१३. --- पोटात \VUgras{जठर} आणि \VUgras{आतडी} असतात. त्यांचा उपयोग अन्न पचवण्यासाठी
होतो.

%%K t02-tit-3 sub
\VUsaut{4.4mm}
\VUpk{10.2pt}{40}{अवयव.}
\VUsaut{3.0mm}

%%K t02-14-1 p
१४. --- आपल्याला चार \VUgras{अवयव} आहेत: दोन \VUgras{हात} आणि दोन \VUgras{पाय}.
हातांनी आपण वस्तू घेतो, धरतो, उचलतो आणि गोळा करतो; पायांनी आपण उभे राहतो,
चालतो आणि धावतो.

%%K t02-15-1 p
१५. --- \VUgras{खांदा} हाताला शरीराशी जोडतो. एक अतिशय बळकट \VUgras{स्नायू},
म्हणजे द्विशिर स्नायू, हाताला हलवतो; आणि \VUgras{कोपर} त्या हाताला
\VUgras{बाहूशी} जोडते. \VUgras{मनगट} हा \VUgras{हाताचा} सांधा आहे, आणि हात
\VUgras{बोटांनी} व \VUgras{नखांनी} संपतो. हातावर आपल्याला \VUgras{शीर} (आणि
रोहिणी) दिसते, जिच्यातून \VUgras{रक्त} वाहते.

%%K t02-16-1 p
१६. --- पायाचे भाग असे आहेत: \VUgras{मांडी}, \VUgras{गुडघा} आणि
\VUgras{गुडघ्यामागची खाच}, \VUgras{पोटरी} --- जिचे \VUgras{मांस}
\VUgras{कातडीने} झाकलेले असते ---, \VUgras{घोटा} आणि \VUgras{पाऊल}. पायाची
\VUgras{हाडे} खूप जाड आणि खूप बळकट असतात. पावलाच्या टोकाला \VUgras{पायाची बोटे}
असतात. उरलेले भाग म्हणजे \VUgras{तळपाय} आणि \VUgras{टाच}.

%%K t02-17-1 p
१७. --- मानवी शरीराचे जवळजवळ पूर्ण वर्णन असे आहे.

%%K t02-17-2 p
\textit{टीप.} --- \textit{«माझे... (डोके वगैरे) दुखते आहे» या वाक्प्रयोगाच्या
साहाय्याने शिक्षक हा सगळा शब्दसंग्रह चांगल्या किंवा वाईट} \VUgras{प्रकृतीच्या}
\textit{दृष्टिकोनातून पुन्हा वापरू शकतील.}

%%K t02-tit-4 sub
\VUsaut{5.0mm}
\VUpk{10.2pt}{40}{व्यायाम.}
\VUsaut{2.4mm}
\VUcentre{10.2pt}{0}{\textit{(एक विद्यार्थी सांगतो आहे.)}}
\VUsaut{3.0mm}

%%K t02-18-1 p
१८. --- लवचिक, चपळ आणि निरोगी राहण्यासाठी आपण आपले पाय आणि आपले हात यांना
व्यायाम दिला पाहिजे. म्हणूनच आपण खेळतो आणि \VUgras{व्यायाम} करतो.

%%K t02-19-1 p
१९. --- आठवड्यात हे दोन्ही व्यायाम शाळेच्या पटांगणात होतात (तक्ता क्रमांक १
पाहा). पण गुरुवारी आणि रविवारी आपण मोठ्या हिरवळीवर जातो, जिथे धावायला आणि
खेळायला जागा असते.

%%K t02-20-1 p
२०. --- आमचे शिक्षक आणि आमचे आईवडील कधीकधी आमच्याबरोबर तिथे येतात; पण व्यायाम
चालवतात ते \VUgras{व्यायामशिक्षक}\textsuperscript{(12)}.

%%K t02-21-1 p
२१. --- हिरवळीवर, प्रवेशाच्या डावीकडे, \VUgras{व्यायामाची
साधने}\textsuperscript{(1)} आहेत; आपण \VUgras{शिडीवर}\textsuperscript{(2)} चढतो,
\VUgras{कड्यांवर}\textsuperscript{(3)} आणि \VUgras{झोक्यावर}\textsuperscript{(4)}
उलटे लटकतो, आणि \VUgras{दोरीच्या शिडीच्या}\textsuperscript{(5)} किंवा
\VUgras{गाठींच्या दोराच्या}\textsuperscript{(6)} टोकापर्यंत हातांनी स्वतःला वर
ओढतो.

%%K t02-22-1 p
२२. --- त्याच वेळी आमचे सोबती \VUgras{लाकडी घोड्यावरून}\textsuperscript{(7)} किंवा
\VUgras{समांतर बारवरून}\textsuperscript{(11)} उड्या मारतात. दुसरे काही
\VUgras{मुष्टियुद्ध करतात}\textsuperscript{(8)} किंवा मुखवटा घालून व
\VUgras{पात्याने} \VUgras{तलवारबाजी करतात}\textsuperscript{(9)}. आणि आणखी काही
\VUgras{वजने उचलतात}\textsuperscript{(10)}.

%%K t02-tit-5 sub
\VUsaut{4.6mm}
\VUpk{10.2pt}{40}{खेळ.}
\VUsaut{3.0mm}

%%K t02-23-1 p
२३. --- व्यायाम करणाऱ्यांच्या भोवती मुले आणि मुली निरनिराळे \VUgras{खेळ} खेळतात.
मुली आपल्या \VUgras{कड्याने}\textsuperscript{(14)} खेळतात;
\VUgras{दोरीवरून}\textsuperscript{(13)} उड्या मारतात, किंवा आपल्या भावांना व
चुलतभावांना \VUgras{क्रोकेच्या}\textsuperscript{(15)} डावासाठी बोलावतात.
प्रतिस्पर्ध्याला वाटते की तो सहज कमानीतून जाईल --- नेमक्या त्याच क्षणी त्याचा
\VUgras{चेंडू} आपल्या \VUgras{लाकडी हातोड्याने} दूर ढकलण्यात त्यांना मौज वाटते!

%%K t02-24-1 p
२४. --- मुलांना अधिक ताकदीचे खेळ आवडतात. \VUgras{ठोकळ्यांचा}\textsuperscript{(16)}
खेळ ते आनंदाने खेळतात, आणि \VUgras{चेंडूने}\textsuperscript{(17)} ते कौशल्याने पाडतात.
\VUgras{भोवरा}\textsuperscript{(18)} आणि \VUgras{बिलबोके}\textsuperscript{(19)}, किंवा
\VUgras{बेडूक-खेळही}\textsuperscript{(20)} ते नाकारत नाहीत. पण त्यांचे सर्वांत आवडते
खेळ म्हणजे \VUgras{गोट्या}\textsuperscript{(21)} आणि \VUgras{भिंतीवरचा
चेंडू}\textsuperscript{(22)}, कारण \VUgras{काचघराची}\textsuperscript{(26)} भिंत हलका
चेंडू परतवण्यासाठी फार सोयीची आहे.

%%K t02-25-1 p
२५. --- लहानगा जॉन आपल्या \VUgras{उंच काठ्यांवरून}\textsuperscript{(24)} चालतो; तो
फार हुशार आहे आणि तोल सांभाळणे त्याला चांगले जमते.

%%K t02-noto-1 noto
\VUnotes{143.2mm}{%
(*) लहान मुलांच्या खेळांना बहुधा पूर्णपणे राष्ट्रीय नावे असतात. ती इदोमध्ये
शक्य तितक्या चांगल्या रीतीने नावे देण्याचा आम्ही प्रयत्न केला: कधी खेळाच्या
कृतीवरूनच, जी त्याचे वैशिष्ट्य असते, तर कधी या किंवा त्या देशातले राष्ट्रीय
नाव घेऊन, जेव्हा ते फारसे अन्याय्य नसते. उदा.: \VUgras{पिंपाचा खेळ} (जर्मन आणि
फ्रेंच) म्हणजेच \VUgras{बेडूक-खेळ} \textsuperscript{(20)} (इंग्रजी), ज्यात
खेळणारे आपल्या गोल चकत्या छिद्रे पाडलेल्या पेटीवर (पिंपावर) तोंड उघडून उभ्या
असलेल्या धातूच्या बेडकाच्या तोंडातून फेकण्याचा प्रयत्न करतात.
\VUgras{खांद्यावरून उडी}\textsuperscript{(30)} आणि \VUgras{पाठीवरून उडी} यांत फक्त
एवढाच फरक आहे: डोक्यावरून उडी मारण्यासाठी खेळणारे आपले हात सोबत्याच्या
खांद्यांवर टेकवतात, तर «\VUgras{पाठीवरून उडीत}» ते हात फक्त त्याच्या पाठीवर
टेकवतात. --- किंवा काही सहभागी वाकून उभे राहतात आणि बाकीचे त्यांच्या पाठीवरून,
खांद्यावर हात टेकवून उड्या मारतात; पहिल्यांनी न वाकता धक्का सोसला पाहिजे आणि
दुसऱ्यांनी तोल न गमावता उडी मारली पाहिजे (तक्ताच पाहा).%
}

%%K t02-25-1 p suite
त्याच्याजवळ काही सोबती त्या काचघरात आणि \VUgras{कुंपणाच्या} मागे
\VUgras{लपाछपी}\textsuperscript{(25)} खेळतात. दुसऱ्यांना \VUgras{कोणी मारले
ओळख}\textsuperscript{(28)} किंवा \VUgras{पिसांचा चेंडू}\textsuperscript{(29)} अधिक
आवडतो, जो एका \VUgras{रॅकेटवरून} दुसऱ्यावर उडत जातो.

%%K t02-26-1 p
२६. --- हिरवळीच्या मधोमध दोन झाडे आहेत. त्यांतल्या एकाच्या बुंध्याशी सात-आठ मुलांनी
\VUgras{खांद्यावरून उडीचा}\textsuperscript{(30)} डाव मांडला आहे (*). हा खेळ सगळ्या
देशांत माहीत नाही; पण \VUgras{पाठीवरून उडी} म्हणजे काय हे सगळ्यांना माहीत आहे ---
आज मात्र मुले तो खेळत नाहीत.

%%K t02-tit-6 sub
\VUsaut{5.0mm}
\VUpk{10.2pt}{40}{मोकळ्या हवेतले खेळ.}
\VUsaut{3.4mm}

%%K t02-27-1 p
२७. \VUgras{आंधळी कोशिंबीरही}\textsuperscript{(31)} फार मजेदार आहे; मुली आणि मुले
आळीपाळीने आपल्या डोळ्यांवर \VUgras{पट्टी} बांधतात आणि जे बेसावध असतात व स्वतःला
पकडू देतात त्यांना धरतात.

%%K t02-28-1 p
२८. --- हिरवळीच्या मधोमध एक अगदी फ्रेंच पण जरा जोरदार खेळ आहे: तो म्हणजे
\VUgras{अस्वल-खेळ}\textsuperscript{(32)}, जो

%%K t02-28-1 p suite
\VUgras{पतंगाच्या}\textsuperscript{(33)} अगदी शांत खेळापेक्षा, आणि
\VUgras{टेनिस}\textsuperscript{(34)} या इंग्रजी खेळापेक्षाही पुष्कळ अधिक गोंगाटाचा
असतो.

%%K t02-29-1 p
२९. --- हा शेवटचा खेळ मोठ्या विद्यार्थ्यांचा आनंद आहे. आमचे शिक्षकसुद्धा तो
आवडीने खेळतात. त्याला चपळाईबरोबर मोठे कौशल्य लागते आणि मला तो फार आवडतो.
\VUgras{झोपाळ्याचे}\textsuperscript{(35)} खेळ मी मुलींसाठी सोडून देतो.

%%K t02-30-1 p
३०. --- आणखी एक इंग्रजी खेळ मला आवडत नाही, जो कदाचित धोक्याचा ठरेल.

%%K t02-30-1b p
मला \VUgras{फुटबॉलविषयी}\textsuperscript{(36)} बोलायचे आहे, जो माझ्या मोठ्या भावाला
आनंद देतो. मला आमचा जुना \VUgras{बार-खेळ}\textsuperscript{(38)} अधिक आवडतो --- तितकाच
आरोग्यदायी आणि खरोखर पुष्कळ कमी रांगडा.

%%K t02-tit-7 sub
\VUsaut{4.6mm}
\VUpk{10.2pt}{40}{सायकल आणि फुगा.}
\VUsaut{3.0mm}

%%K t02-31-1 p
३१. --- \VUgras{सायकलवरून}\textsuperscript{(37)} हिरवळीभोवती फेरी मारायलाही मला
आवडते.

%%K t02-32-1 p
३२. --- पुढच्या वर्षी मी \VUgras{घोडेस्वारी}\textsuperscript{(39)} शिकणार आहे. डौलाने
चालणारा किंवा दुडक्या चालीने जाणारा, आणि सरपटत अडथळ्यांवरून उडी मारणारा घोडा
किती सुंदर दिसतो!

%%K t02-33-1 p
३३. --- उन्हाळ्यात आपण \VUgras{पोहायला}\textsuperscript{(40)} शिकतो. नदीच्या स्वच्छ
आणि थंड पाण्यात आपण आनंदाने बुडी मारतो आणि डुंबतो. कधीकधी शेवटी आपण कागदाचा
\VUgras{फुगा}\textsuperscript{(41)} सोडतो, जो गरम हवेने भरताच डौलाने आकाशात चढतो.

%%K t02-34-1 p
३४. --- आणखीही पुष्कळ खेळ आहेत, पण ते मोजून मला संपणार नाही; आणि या थोडक्या
वर्णनावरून कोणाच्याही लक्षात येईल की गुरुवारी आमच्या सुंदर हिरवळीवर कोणाला
फारसा कंटाळा येत नाही.

%%K t02-34-2 p
टीप. --- \textit{शिक्षक हे प्रकरण सहज पूर्ण करू शकतील; सर्वांत महत्त्वाच्या
प्रत्येक खेळाचे अधिक तपशीलवार वर्णन करून.}
